\documentclass[a4paper,10pt]{article}

% ---------- CONFIGURAÇÃO ----------
\usepackage[
    a4paper,
    top=0.65cm,
    bottom=0.65cm,
    left=1.25cm,
    right=1.25cm
]{geometry}

\usepackage[T1]{fontenc}
\usepackage[utf8]{inputenc}
\usepackage[portuguese,provide=*]{babel}
\IfFileExists{newtxtext.sty}{
  \usepackage{newtxtext}
  \usepackage{newtxmath}
}{
  \usepackage{mathptmx}
}
\usepackage{amssymb}
\usepackage{microtype}
\usepackage{enumitem}
\usepackage{hyperref}
\usepackage{titlesec}
\usepackage{xcolor}

\definecolor{linkcolor}{RGB}{0, 0, 160}

\hypersetup{
    colorlinks=true,
    urlcolor=linkcolor,
    linkcolor=linkcolor,
    pdftitle={Israel Souza Ferreira - Currículo},
    pdfauthor={Israel Souza Ferreira}
}

\urlstyle{same}

\setlength{\parindent}{0pt}
\setlength{\parskip}{0pt}

\pagestyle{empty}
\raggedbottom
\emergencystretch=1em

% ---------- LISTAS COMPACTAS ----------
\setlist[itemize]{
    label=\textbullet,
    leftmargin=1.2em,
    itemsep=0.25pt,
    topsep=1pt,
    parsep=0pt,
    partopsep=0pt
}

% ---------- SEÇÕES ----------
\titleformat{\section}
  {\normalsize\bfseries}
  {}
  {0em}
  {\MakeUppercase}
  [\vspace{-1.1ex}\rule{\linewidth}{0.35pt}]

\titlespacing{\section}{0pt}{1.05ex}{0.4ex}

% #1 Organização
% #2 Local
% #3 Cargo ou formação
% #4 Período
\newcommand{\entrada}[4]{%
    \noindent
    \textbf{#1}
    \hfill
    \textbf{#2}\\[-1pt]
    \textit{#3}
    \hfill
    \textit{#4}\par
}

\begin{document}

% ---------- CABEÇALHO ----------
\begin{center}

    {\LARGE\textbf{Israel Souza Ferreira}}\\[1pt]

    {\normalsize
        \textbf{Cientista de Dados
        \enspace $|$ \enspace
        Engenheiro de Machine Learning}
    }\\[2pt]

    {\small
        Fortaleza, CE
        \enspace $|$ \enspace
        Remoto
        \enspace $|$ \enspace
        +55 85 98433-2790
        \enspace $|$ \enspace
        \href{mailto:israel.souza.dev@gmail.com}
        {israel.souza.dev@gmail.com}
        \\[1pt]
        \href{https://www.linkedin.com/in/isrreal/}
        {linkedin.com/in/isrreal}
        \enspace $|$ \enspace
        \href{https://github.com/isrreal}
        {github.com/isrreal}
        \enspace $|$ \enspace
        \href{https://israelsouza.speculummaius.com.br/}
        {Portfólio: israelsouza.speculummaius.com.br}
        \enspace $|$ \enspace
        \href{http://lattes.cnpq.br/6540409845091993}
        {Lattes}
    }

\end{center}

\vspace{-3pt}

% ---------- PERFIL ----------
\section*{Perfil}

Cientista de dados e engenheiro de machine learning com experiência em
\textbf{pesquisa aplicada} e entrega ponta a ponta: experimentação, visão
computacional, APIs, dados e implantação CPU-only. Transformo hipóteses em
decisões mensuráveis com baselines, ablações, benchmarks e testes automatizados.

% ---------- EXPERIÊNCIA ----------
\section*{Experiência Profissional}

\entrada
    {Tieta Artificial Intelligence}
    {Remoto, Brasil}
    {Pesquisador de P\&D --- Programa de Formação em Pesquisa e Inovação}
    {Jul 2026 - Presente}

\begin{itemize}

    \item Desenvolvo um pipeline de \textbf{OCR e segmentação de documentos}
    que transforma provas escaneadas em questões estruturadas para análise
    educacional.

\end{itemize}

\entrada
    {Tieta Artificial Intelligence}
    {Remoto, Brasil}
    {Cientista de Dados / Pesquisador de P\&D --- Bolsista CNPq RHAE}
    {Abr 2025 - Jun 2026}

\begin{itemize}

    \item Modelei parâmetros psicométricos com \textbf{Teoria de Resposta ao
    Item} e comparei abordagens por avaliações pareadas, ablações, correlação
    com referências e custo de inferência.

    \item Construí pipelines de extração de sinais com LLMs orquestrados por
    \textbf{DSPy} e avaliei arquiteturas de regressão multitarefa sobre dados
    textuais e visuais em experimentos reprodutíveis.

\end{itemize}

\entrada
    {Consultoria autônoma --- Instituição de ensino}
    {Remoto, Brasil}
    {Desenvolvedor Full Stack e Engenheiro de Machine Learning}
    {2026 - Presente}

\begin{itemize}

    \item Desenvolvi ponta a ponta o \href
    {https://github.com/isrreal/face-clock-evoluir-public}
    {\textbf{Face Clock Evoluir}}, sistema contratado de ponto facial em
    homologação final, cobrindo modelos, backend, banco, segurança, testes e
    implantação.

    \item Coloquei reconhecimento 1:N em operação CPU-only com YuNet + SFace,
    ONNX Runtime e PostgreSQL/pgvector: \textbf{p95 de 74~ms} com uma
    identificação e \textbf{574 testes automatizados} na entrega atual.

    \item Recalibrei o limiar open-set, reduzindo desconhecidos aceitos de
    \textbf{45,4\% para 0,2\%} na base pública; reduzi \textbf{5 para 3 frames}
    após medir acurácia idêntica de 1 a 8 e corrigi uma corrida que causava
    \textbf{24/24 falhas} em oito rajadas simultâneas.

    \item Um benchmark de busca vetorial mediu \textbf{0,076~ms em 1.000
    vetores}; como o HNSW obteve recall@1 de 0,915, mantive a busca exata e
    transacional no pgvector em vez de adotar FAISS prematuramente.

    \item Comparei cinco abordagens para triagem de atestados; a solução
    EfficientNet + Regressão Logística alcançou \textbf{87,5\% de acurácia} e
    sete falsos positivos a menos, preservando o recall.

    \item Adequei a entrega ao servidor sem GPU e com 8~GB de RAM: imagens de
    produção de \textbf{10,59 para 1,11~GB} (9,5$\times$ menores), com zero
    decisões do classificador alteradas na migração para ONNX Runtime.

\end{itemize}

\entrada
    {Universidade Federal do Ceará --- Campus Quixadá}
    {Presencial, Brasil}
    {Monitor de Algoritmos, Cálculo e Pré-Cálculo}
    {2022 - 2024}

\begin{itemize}

    \item Apoiei turmas em estruturas de dados, algoritmos, complexidade e
    Cálculo, traduzindo conteúdo quantitativo para diferentes níveis.

\end{itemize}

% ---------- PROJETOS ----------
\section*{Projetos Selecionados}

\begin{itemize}

    \item
    \textbf{
        \href
        {https://github.com/isrreal/Triple-Roman-Domination-in-graphs}
        {Triple Roman Domination in Graphs (TCC):}
    }
    primeiras meta-heurísticas GA e ACO da literatura para o problema e
    correção da formulação PLI publicada; avaliação sobre \textbf{362 grafos}.

    \item
    \textbf{
        \href
        {https://github.com/isrreal/AuxilioEmergencialQueries}
        {Auxílio Emergencial --- Engenharia de Dados:}
    }
    ingestão em blocos com \texttt{COPY} binário sobre \textbf{31,6~GB} de
    dados públicos; instrumentação ligou o crescimento de memória aos IDs de
    deduplicação retidos (\textbf{r = 0,99}).

\end{itemize}

% ---------- HABILIDADES ----------
\section*{Habilidades Técnicas}

\textbf{Linguagens e consultas:}
Python, SQL, C++ e Bash.\\[1pt]

\textbf{ML, visão e experimentação:}
PyTorch, Scikit-learn, OpenCV, YuNet, SFace, ONNX Runtime, DSPy, Optuna,
validação cruzada, calibração, ablação, testes de hipótese e busca vetorial.\\[1pt]

\textbf{Dados e produção:}
PostgreSQL, pgvector, FAISS, Pandas, NumPy, SciPy, FastAPI, SQLAlchemy, Alembic,
Docker, GitHub Actions, pytest, MLflow e Streamlit.

% ---------- FORMAÇÃO ----------
\section*{Formação Acadêmica}

\entrada
    {Universidade Federal do Ceará --- Campus Quixadá}
    {Brasil}
    {Bacharelado em Ciência da Computação}
    {Conclusão: Jul 2026}

% ---------- IDIOMAS ----------
\section*{Idiomas}

\textbf{Português:} nativo
\enspace $|$ \enspace
\textbf{Inglês:} leitura técnica fluente de artigos e documentações

\end{document}
